\scanpage{184}

UNIVERSITE DES SCIENCES\\
ET TECHNIQUES DU LANGUEDOC
\hfill SESSION DE JUIN 1981

\begin{center}
D.E.U.G. A 1ère Année - Maths Appliquées\\
(Séries I et II)

Epreuve d'Analyse - Durée 3 heures

I
\end{center}

On donne l'équation différentielle
\[
xyy'+2y^2=\frac{1}{\unsure{x}^3(x+1)^2}
\]

1 - Déterminer la solution de cette équation définie sur l'intervalle $]1,+\infty[$
et qui prend la valeur $-\frac{1}{4\sqrt3}$ pour $x=2$.

On pourra poser $y^2=z$.

2 - Soit A un nombre réel ; $A>1$. Calculer l'intégrale
\[
\mathrm F(A)=\int_1^A\frac{1}{x^2}\sqrt{\frac{x-1}{x+1}}\,\mathrm dx
\]
\unsure{On pourra poser} $x\unsure{=}\frac1u$ \unsure{et utiliser} la relation
\[
\sqrt{\frac{1-u}{1+u}}=\frac{1-u}{\sqrt{1-u^2}}
\]

3 - Etudier la limite de $\mathrm F(A)$ lorsque A tend vers $+\infty$.

\begin{center}II\end{center}

Etude de la courbe définie paramétriquement par les équations
\[
x=\frac{t(3t^2-2)}{1-t^2}
\qquad
y=\frac{t^4}{1-t^2}
\]

1 - Montrer qu'il suffit d'étudier l'arc obtenu pour $t\in[0,+\infty[$,
l'ensemble de la courbe pouvant s'en déduire par une symétrie.

2 - Calculer les dérivées $\frac{\mathrm dx}{\mathrm dt}$, $\frac{\mathrm dy}{\mathrm dt}$.
Vérifier que ces deux dérivées s'annulent pour $t=\sqrt2$.
Etablir le tableau donnant les variations de x et de y en fonction de t.

3 - Etude de la courbe au voisinage de l'origine.

4 - Etude de la \unsure{branche} infinie obtenue lorsque t tend vers 1.

5 - Etude de la branche infinie obtenue lorsque t tend vers $+\infty$.

On pourra chercher les parties entières de x et de y considérées
comme fractions rationnelles de t et en déduire une parabole asymptotique à la
courbe.

6 - Etude de la courbe au voisinage du point $t=\sqrt2$.

7 - Dessin de la courbe entière.

\begin{center}\rule{1.5cm}{.4pt}\end{center}

\margin{184}
